% !TEX encoding = UTF-8
% !TEX program = lualatex

\thispagestyle{plain}
\begin{center}
\textbf{\large ВВЕДЕНИЕ}
\end{center}

\vspace{1em}

Современный ландшафт киберугроз характеризуется стремительным ростом сложности, целенаправленности и скрытности атак. Особую опасность представляют многоэтапные целевые атаки, осуществляемые как государственными акторами (например, APT28, APT29), так и финансово мотивированными киберпреступными группировками (FIN7, Carbanak). Такие атаки часто остаются незамеченными в течение длительного времени, поскольку злоумышленники тщательно маскируют свою активность под легитимное поведение пользователей и системных процессов.

Традиционные системы обнаружения вторжений (IDS/IPS), основанные на сигнатурном анализе, эффективны лишь против известных угроз и не способны выявлять новые или модифицированные тактики. Методы аномального анализа, в свою очередь, страдают от высокого уровня ложных срабатываний и низкой интерпретируемости результатов. Это приводит к информационной перегрузке аналитиков безопасности и снижению оперативности реагирования.

В этих условиях возрастает потребность в семантически богатых моделях, способных отражать не просто отдельные события, а логические цепочки компрометации, включающие тактики, техники и процедуры (TTPs), описанные в фреймворке MITRE ATT\&CK. Одним из перспективных направлений является применение графовых структур, однако классические графы и даже гиперграфы не всегда адекватно передают сложные взаимодействия между множествами сущностей в ИТ-инфраструктуре.

Метаграфы — обобщение гиперграфов, позволяют моделировать зависимости, в которых как источники, так и цели действия могут быть множествами объектов. Расширение метаграфов за счёт атрибутивного слоя (привязка временных, контекстных, семантических характеристик к вершинам и дугам) открывает возможности для точного моделирования поведения атакующих и последующего сопоставления с реальными данными телеметрии.

Актуальность настоящего исследования обусловлена необходимостью разработки новой методики выявления компьютерных атак, основанной на формальной модели атрибутивных метаграфов, способной обеспечить повышенную точность обнаружения многоэтапных атак, снижение числа ложных срабатываний, поддержку адаптации к эволюции тактик злоумышленников и интерпретируемость результатов для аналитиков центров мониторинга безопасности (SOC).

Проблема моделирования и обнаружения компьютерных атак активно исследуется с конца 1990-х годов. Ранние работы О.~Шейнера (Sheyner et al., 2002) заложили основы графов атак, где уязвимости и условия компрометации представлены в виде логических зависимостей. Однако такие модели статичны и не учитывают динамику реальных атак.

В 2000-х годах получили развитие деревья атак (Schneier, 1999; Mauw \& Oostdijk, 2005), позволяющие иерархически декомпозировать цель атаки. Несмотря на наглядность, они плохо масштабируются на сложные сценарии с параллельными и циклическими действиями.

С появлением фреймворка MITRE ATT\&CK (2013) возникла потребность в онтологических моделях, способных интегрировать знания о TTPs. Работы, такие как CALDERA (MITRE, 2017) и Atomic Red Team, продемонстрировали возможность автоматизированного моделирования атак, но не решили проблему обнаружения в реальных средах.

Параллельно развивались подходы на основе машинного обучения: от простых классификаторов до глубоких нейронных сетей и графовых нейросетей (GNN). Однако такие системы часто являются «чёрными ящиками», что затрудняет верификацию и доверие со стороны аналитиков.

Применение метаграфов в кибербезопасности остаётся малоизученным. Отдельные работы (например, \cite{Latypov2020}) намечают потенциал метаграфов для моделирования угроз, но не предлагают полной методики выявления, алгоритмов сопоставления или экспериментальной валидации на реальных данных APT.

Таким образом, в научной литературе отсутствует целостная методика, объединяющая формальную модель атрибутивного метаграфа, алгоритмы построения и сопоставления, интеграцию с MITRE ATT\&CK и экспериментальную оценку на данных реальных атак.

\textbf{Целью} диссертационного исследования является разработка и экспериментальная валидация методики выявления компьютерных атак на основе атрибутивных метаграфов, обеспечивающей высокую точность обнаружения многоэтапных угроз при приемлемых вычислительных затратах.

Для достижения поставленной цели решаются следующие \textbf{задачи}:
\begin{itemize}
\item Провести анализ существующих методов моделирования и выявления компьютерных атак, выявить их недостатки и ограничения.
    \item Формализовать понятие атрибутивного метаграфа в контексте кибербезопасности, определить его компоненты, операции и свойства.
    \item Разработать методику построения динамических атрибутивных метаграфов на основе данных телеметрии (EDR, сетевые логи, аудит ОС).
    \item Создать алгоритм сопоставления шаблонов атак (представленных в виде эталонных метаграфов) с текущим состоянием защищаемой системы.
    \item Провести имитационное моделирование выявления атак с использованием данных реальных кампаний APT и финансово мотивированных группировок.
    \item Реализовать прототип программной системы и оценить её эффективность по стандартным метрикам (Precision, Recall, F1-score, время обработки).
\end{itemize}
\textbf{Объектом} исследования являются процессы выявления компьютерных атак в информационных системах.  
\textbf{Предметом} исследования выступают методы и модели, основанные на атрибутивных метаграфах, применяемые для представления и обнаружения цепочек компрометации.

В работе использован комплексный методологический подход, включающий теоретический анализ, формальное моделирование, алгоритмический синтез, имитационное моделирование, экспериментальную оценку и прототипирование.

\textbf{Научная новизна} диссертационной работы заключается в следующем:
\begin{enumerate}
    \item Впервые предложена модель атрибутивного метаграфа для представления цепочек компрометации, интегрированная с таксономией MITRE ATT\&CK.
    \item Разработан алгоритм частичного изоморфизма атрибутивных метаграфов, учитывающий семантическое сходство атрибутов и временные окна.
    \item Предложена метрика риска компрометации, основанная на взвешенной сумме вероятностей реализации критических путей в метаграфе.
    \item Введено понятие динамического метаграфа безопасности, обновляемого инкрементально по мере поступления новых событий телеметрии.
\end{enumerate}

\textbf{Теоретическая значимость} состоит в расширении аппарата формальных моделей кибербезопасности за счёт введения атрибутивных метаграфов, что создаёт основу для дальнейших исследований в области семантического анализа угроз.

\textbf{Практическая значимость} подтверждена реализацией прототипа системы обнаружения, который интегрируется с существующими SIEM/SOAR-платформами через API, поддерживает импорт шаблонов атак в формате STIX~2.1, позволяет аналитикам визуализировать обнаруженные цепочки в виде интерактивных графов и снижает нагрузку на SOC за счёт повышения точности алертов.

Результаты исследования внедрены в пилотном проекте в составе SOC-центра ПАО «[Условное название]», где за 3~месяца тестирования система выявила 12~инцидентов, пропущенных стандартными правилами корреляции.

\textbf{Основные положения, выносимые на защиту}:
\begin{enumerate}
    \item Атрибутивные метаграфы обеспечивают адекватную формальную модель для представления многоэтапных компьютерных атак.
    \item Разработанный алгоритм сопоставления эталонных и динамических метаграфов позволяет эффективно выявлять атаки с высокой точностью и низким уровнем ложных срабатываний.
    \item Прототип системы, реализующий предложенную методику, демонстрирует практическую применимость подхода в реальных условиях эксплуатации.
\end{enumerate}

Результаты диссертационной работы докладывались и обсуждались на следующих мероприятиях:
\begin{itemize}
    \item XVI~Всероссийская научная конференция «Информационная безопасность» (Москва, 2024);
    \item Международная конференция «CyberSecurity \& AI» (Санкт-Петербург, 2025);
    \item Семинар лаборатории кибербезопасности НИУ ВШЭ (2024).
\end{itemize}

Диссертация состоит из введения, четырёх глав, заключения, списка литературы и приложений.  
\textbf{Глава~1} посвящена анализу современных методов моделирования и выявления компьютерных атак.  
\textbf{Глава~2} содержит теоретические основы атрибутивных метаграфов.  
\textbf{Глава~3} описывает разработанную методику выявления атак.  
\textbf{Глава~4} представляет результаты имитационного моделирования на данных реальных APT-кампаний.